\subsubsection*{Assignment 3 -- Song Profiling}
L'Assignment 3 richiedeva di definire categorie di canzoni,
basandosi sulle informazioni riguardanti i lyrics e/o altre informazioni.
La classificazione è stata basata esclusivamente sul testo, applicando
un modello di topic modeling.
Il processo è stato articolato nei seguenti passaggi:

\begin{enumerate}
    \item \textbf{Preprocessing}: le lyrics sono state lemmatizzate con
    \texttt{spaCy}, distinguendo automaticamente tra inglese e italiano,
    e filtrando stopwords generali e specifiche del dataset.
    
    \item \textbf{Rappresentazione testuale}: costruzione di una matrice
    documento-termine tramite TF-IDF.
    
    \item \textbf{Modellazione dei topic}: addestramento di un modello
    NMF con 8 topic, ottenendo la matrice termine-topic e
    documento-topic.
    
    \item \textbf{Assegnazione delle categorie}: ogni canzone è stata
    associata al topic dominante.
    Basandosi sulle parole più presenti in ogni topic, ad ognuno di essi
    sono stati assegnati nomi interpretabili.
\end{enumerate}

Il dataset finale contiene quindi, per ogni canzone con lyrics
disponibili, un attributo \texttt{category} che ne descrive il genere
tematico principale.